\chapter{Practical System Design}

\section{Framing ML as a Systems Problem}

Machine learning is often presented as a problem of choosing a model and minimizing a loss. In practice, however, useful machine learning systems are much larger than the model itself. They involve data collection, labeling, storage, training infrastructure, evaluation, deployment, monitoring, and continual maintenance.

\medskip

\noindent
\textbf{Core viewpoint.}
A machine learning system is not just a function
\[
x \mapsto \hat{y},
\]
but a pipeline that connects:
\begin{itemize}
    \item raw data,
    \item human decisions and labels,
    \item training and validation workflows,
    \item deployed inference services,
    \item monitoring and feedback loops.
\end{itemize}

\medskip

\noindent
\textbf{Why this matters.}
A theoretically strong model can still fail in practice if:
\begin{itemize}
    \item the data pipeline is unreliable,
    \item the evaluation setup is flawed,
    \item the system violates latency or memory constraints,
    \item deployment conditions differ from training conditions,
    \item failures are not detected after launch.
\end{itemize}

\medskip

\noindent
Thus practical machine learning must be framed as a \emph{systems problem}, not merely an optimization problem.

\medskip

\noindent
\textbf{Main system components.}
\begin{itemize}
    \item \textbf{Data infrastructure:} collection, storage, cleaning, labeling, versioning
    \item \textbf{Training infrastructure:} preprocessing, batching, training, validation, checkpointing
    \item \textbf{Experiment infrastructure:} metrics, logs, hyperparameters, comparisons
    \item \textbf{Serving infrastructure:} model packaging, inference, latency, scaling
    \item \textbf{Monitoring infrastructure:} drift detection, error analysis, feedback collection
\end{itemize}

\medskip

\noindent
\textbf{Insight.}
In real-world ML, system quality is determined by the weakest link in the pipeline, not only by the sophistication of the learning algorithm.

\section{Data Collection, Cleaning, and Labeling}

Data is the foundation of a machine learning system. In many applications, the quality of the dataset matters more than small differences between model architectures.

\subsection{Data Collection}

The first step is to gather data relevant to the task.

\medskip

\noindent
\textbf{Typical sources.}
\begin{itemize}
    \item sensors and devices,
    \item user interactions,
    \item logs and transaction records,
    \item public datasets,
    \item manually curated examples.
\end{itemize}

\medskip

\noindent
\textbf{Key concerns.}
\begin{itemize}
    \item Does the data represent the target population?
    \item Does it cover the range of situations the model will face?
    \item Are rare but important cases included?
\end{itemize}

\medskip

\noindent
A common failure is \emph{sampling mismatch}: the collected data is easier, cleaner, or less diverse than real deployment data.

\subsection{Cleaning and Preprocessing}

Raw data is rarely ready for immediate use. Practical systems require systematic preprocessing.

\medskip

\noindent
\textbf{Examples.}
\begin{itemize}
    \item removing duplicates,
    \item handling missing values,
    \item correcting corrupted records,
    \item normalizing units and formats,
    \item standardizing text, images, or numerical features.
\end{itemize}

\medskip

\noindent
\textbf{Feature preprocessing.}
Depending on the modality, one may apply:
\begin{itemize}
    \item normalization or scaling,
    \item tokenization,
    \item image resizing and augmentation,
    \item categorical encoding,
    \item filtering and denoising.
\end{itemize}

\medskip

\noindent
\textbf{Pipeline principle.}
Preprocessing should be automated and reproducible. Ad hoc manual transformations create hidden inconsistencies between training and deployment.

\subsection{Labeling}

Supervised learning depends on labels, and labeling quality is often a major bottleneck.

\medskip

\noindent
\textbf{Label sources.}
\begin{itemize}
    \item expert annotation,
    \item crowd annotation,
    \item user-provided outcomes,
    \item weak or heuristic labels,
    \item labels inferred from downstream events.
\end{itemize}

\medskip

\noindent
\textbf{Common labeling issues.}
\begin{itemize}
    \item inconsistent annotation standards,
    \item ambiguous examples,
    \item class imbalance,
    \item noisy labels,
    \item delayed or missing ground truth.
\end{itemize}

\medskip

\noindent
Because models learn the structure of labels, poor labeling policy often leads to systematic model failure.

\subsection{Train, Validation, and Test Infrastructure}

A robust system requires careful separation of datasets.

\medskip

\noindent
\textbf{Training set.}
Used to fit model parameters.

\medskip

\noindent
\textbf{Validation set.}
Used to tune hyperparameters, choose architectures, and make model-selection decisions.

\medskip

\noindent
\textbf{Test set.}
Used only for final evaluation. It should remain isolated from repeated experimentation.

\medskip

\noindent
\textbf{Why the split matters.}
If the validation or test data leaks into training decisions, performance estimates become overly optimistic.

\medskip

\noindent
\textbf{Practical infrastructure concerns.}
\begin{itemize}
    \item prevent duplicate or near-duplicate examples across splits,
    \item ensure time-based splits when future prediction is the goal,
    \item preserve group boundaries when examples are correlated,
    \item version the exact train/val/test partitions.
\end{itemize}

\medskip

\noindent
A train/validation/test split is not only a statistical idea but also an engineering artifact that must be stored, reproduced, and audited.

\subsection{Data Pipelines}

A \textbf{data pipeline} transforms raw data into model-ready inputs and labels.

\medskip

\noindent
A concrete pipeline may include:
\begin{enumerate}
    \item ingest raw records,
    \item validate schema and basic integrity,
    \item clean and standardize values,
    \item join multiple sources,
    \item generate features or encodings,
    \item write processed data to storage for training or serving.
\end{enumerate}

\medskip

\noindent
\textbf{Batch versus streaming pipelines.}
\begin{itemize}
    \item \textbf{Batch pipelines} process data periodically in large jobs
    \item \textbf{Streaming pipelines} process examples continuously as they arrive
\end{itemize}

\medskip

\noindent
The correct choice depends on the application. Offline recommendation retraining may tolerate batch updates, while fraud detection may require continuous ingestion.

\medskip

\noindent
\textbf{Insight.}
Reliable data pipelines are essential because every downstream stage depends on their correctness.

\section{Training Pipelines and Experiment Management}

Training a model in practice involves much more than running gradient descent. One needs infrastructure for repeatable experiments, efficient compute use, and trustworthy comparisons.

\subsection{Training Pipelines}

A training pipeline coordinates data loading, optimization, evaluation, and checkpointing.

\medskip

\noindent
\textbf{Typical stages.}
\begin{enumerate}
    \item read processed training data,
    \item construct batches,
    \item compute forward pass,
    \item evaluate loss,
    \item compute gradients,
    \item update parameters,
    \item periodically evaluate on validation data,
    \item save checkpoints and logs.
\end{enumerate}

\medskip

\noindent
\textbf{Operational requirements.}
\begin{itemize}
    \item training must resume after interruption,
    \item failed jobs should be recoverable,
    \item metrics should be stored continuously,
    \item model artifacts should be versioned.
\end{itemize}

\subsection{Experiment Tracking}

Modern model development requires many runs with different hyperparameters, architectures, and preprocessing choices.

\medskip

\noindent
\textbf{Experiment tracking should record:}
\begin{itemize}
    \item code version,
    \item dataset version,
    \item train/val/test split,
    \item hyperparameters,
    \item optimizer settings,
    \item metrics and curves,
    \item final model artifact.
\end{itemize}

\medskip

\noindent
Without systematic experiment tracking, it becomes difficult to answer basic questions such as:
\begin{itemize}
    \item Which change improved performance?
    \item Which dataset version was used?
    \item Can this result be reproduced?
\end{itemize}

\medskip

\noindent
\textbf{Practical principle.}
Every important result should be traceable to a specific configuration of data, code, and compute.

\subsection{Hyperparameter Search}

Training quality often depends strongly on hyperparameters such as:
\begin{itemize}
    \item learning rate,
    \item regularization strength,
    \item batch size,
    \item architecture depth or width,
    \item augmentation policy.
\end{itemize}

\medskip

\noindent
Common strategies include:
\begin{itemize}
    \item grid search,
    \item random search,
    \item adaptive methods such as Bayesian optimization,
    \item multi-fidelity methods with early stopping.
\end{itemize}

\medskip

\noindent
Hyperparameter tuning is itself a systems task because it requires scheduler logic, job parallelism, budget control, and fair comparison criteria.

\subsection{Checkpointing and Recovery}

Long training jobs require checkpointing.

\medskip

\noindent
\textbf{Why checkpointing matters.}
\begin{itemize}
    \item it prevents loss of progress,
    \item it enables later evaluation or fine-tuning,
    \item it allows rollback to earlier model states,
    \item it supports comparison of intermediate learning stages.
\end{itemize}

\medskip

\noindent
A checkpoint commonly stores:
\begin{itemize}
    \item model parameters,
    \item optimizer state,
    \item scheduler state,
    \item training step or epoch count,
    \item metadata about the run.
\end{itemize}

\subsection{Distributed Training}

Large-scale models may exceed the capacity of a single device.

\medskip

\noindent
\textbf{Two major patterns.}
\begin{itemize}
    \item \textbf{Data parallelism:} replicate the model across devices and split batches
    \item \textbf{Model parallelism:} split the model itself across devices
\end{itemize}

\medskip

\noindent
These choices affect throughput, communication cost, memory use, and engineering complexity.

\medskip

\noindent
\textbf{Insight.}
Efficient training pipelines make experimentation faster, more reliable, and easier to interpret.

\section{Serving, Inference, and Latency Constraints}

A trained model is only useful if it can be deployed in a form compatible with real-world constraints.

\subsection{Serving}

\textbf{Serving} is the process of making model predictions available to downstream users or applications.

\medskip

\noindent
Deployment may occur in several ways:
\begin{itemize}
    \item \textbf{online serving:} predictions are generated on request,
    \item \textbf{batch inference:} predictions are computed periodically for many items,
    \item \textbf{on-device inference:} the model runs locally on edge hardware.
\end{itemize}

\medskip

\noindent
Each setting imposes different requirements on memory, throughput, latency, and reliability.

\subsection{Inference Constraints}

At training time one often optimizes predictive accuracy. At serving time, additional constraints become central:
\begin{itemize}
    \item inference latency,
    \item throughput,
    \item memory footprint,
    \item power consumption,
    \item hardware compatibility,
    \item cost per prediction.
\end{itemize}

\medskip

\noindent
A model with slightly better accuracy may be impractical if it is too slow or too expensive to serve.

\subsection{Latency}

\textbf{Inference latency} is the time required to produce a prediction after an input arrives.

\medskip

\noindent
Latency matters especially in:
\begin{itemize}
    \item interactive systems,
    \item ranking and recommendation,
    \item autonomous systems,
    \item real-time anomaly or fraud detection.
\end{itemize}

\medskip

\noindent
\textbf{Sources of latency.}
\begin{itemize}
    \item preprocessing time,
    \item feature retrieval,
    \item network communication,
    \item model execution,
    \item postprocessing and business logic.
\end{itemize}

\medskip

\noindent
Thus latency is a property of the entire system, not only of the model.

\subsection{Deployment Optimizations}

Common ways to reduce inference cost include:
\begin{itemize}
    \item model pruning,
    \item quantization,
    \item distillation,
    \item caching frequent predictions,
    \item batching requests when possible,
    \item specialized hardware acceleration.
\end{itemize}

\medskip

\noindent
These methods trade off accuracy, complexity, and operational efficiency.

\subsection{Training--Serving Skew}

A major practical problem is \textbf{training--serving skew}: the model is trained on one data representation but receives a different one in deployment.

\medskip

\noindent
This can happen if:
\begin{itemize}
    \item feature definitions differ between pipelines,
    \item missing values are handled differently,
    \item preprocessing code is duplicated inconsistently,
    \item online feature stores are out of sync with offline data.
\end{itemize}

\medskip

\noindent
A robust system should minimize the gap between training and serving pipelines.

\medskip

\noindent
\textbf{Insight.}
Deployment success depends not only on predictive quality, but on whether the model fits operational constraints and receives the inputs it was designed for.

\section{Monitoring Drift, Failures, and Feedback Loops}

Deployment is not the end of the machine learning lifecycle. Models must be monitored continuously after release.

\subsection{Why Monitoring Is Necessary}

Real environments change over time:
\begin{itemize}
    \item user behavior changes,
    \item sensors degrade,
    \item markets shift,
    \item interfaces change,
    \item upstream systems evolve.
\end{itemize}

\medskip

\noindent
A model that performed well at launch may deteriorate later if the data distribution changes.

\subsection{Drift}

A central monitoring concern is \textbf{drift}.

\medskip

\noindent
\textbf{Covariate drift.}
The distribution of inputs changes:
\[
P_{\mathrm{train}}(x) \neq P_{\mathrm{deploy}}(x).
\]

\medskip

\noindent
\textbf{Label or concept drift.}
The relationship between inputs and targets changes:
\[
P_{\mathrm{train}}(y \mid x) \neq P_{\mathrm{deploy}}(y \mid x).
\]

\medskip

\noindent
Both types of drift can degrade performance.

\subsection{Operational Monitoring}

A deployed system should monitor:
\begin{itemize}
    \item input distributions,
    \item prediction distributions,
    \item latency and throughput,
    \item system errors and timeouts,
    \item downstream task metrics when ground truth becomes available.
\end{itemize}

\medskip

\noindent
\textbf{Examples of warning signs.}
\begin{itemize}
    \item sudden change in average predicted score,
    \item missing or malformed input features,
    \item growing disagreement with human review,
    \item increased response time,
    \item drop in business or task outcomes.
\end{itemize}

\subsection{Failures and Silent Errors}

Some ML failures are obvious, but many are silent.

\medskip

\noindent
\textbf{Examples of silent failures.}
\begin{itemize}
    \item shifted label semantics,
    \item broken feature extraction,
    \item incorrect join logic in pipelines,
    \item stale cached features,
    \item model serving the wrong version.
\end{itemize}

\medskip

\noindent
Because such errors may not crash the system, explicit monitoring is essential.

\subsection{Feedback Loops}

ML systems often influence the very data they later receive.

\medskip

\noindent
\textbf{Examples.}
\begin{itemize}
    \item a recommendation system changes what users see and click,
    \item a ranking model changes exposure of content,
    \item a fraud model changes which transactions are investigated,
    \item a predictive system influences human decisions and labels.
\end{itemize}

\medskip

\noindent
This creates \textbf{feedback loops}, in which deployment changes future training data.

\medskip

\noindent
Feedback loops can create:
\begin{itemize}
    \item selection bias,
    \item reinforcing errors,
    \item reduced exploration,
    \item instability over long horizons.
\end{itemize}

\medskip

\noindent
Therefore monitoring should include not only raw performance, but also how the model changes the environment in which it operates.

\section{Engineering Tradeoffs in Real-World ML}

Practical machine learning always involves tradeoffs rather than a single objective.

\subsection{Accuracy versus Simplicity}

A larger model may achieve slightly better validation accuracy, but a smaller model may be preferable if it is:
\begin{itemize}
    \item easier to debug,
    \item easier to deploy,
    \item faster to serve,
    \item more stable across retraining runs.
\end{itemize}

\subsection{Latency versus Quality}

In real-time applications, one may sacrifice some predictive accuracy to satisfy strict latency budgets.

\medskip

\noindent
For example, a two-stage architecture may use:
\begin{itemize}
    \item a fast model for candidate filtering,
    \item a more expensive model for final reranking.
\end{itemize}

\subsection{Freshness versus Stability}

Frequent retraining may adapt quickly to new data, but it can also introduce instability and make debugging harder.

\medskip

\noindent
Less frequent retraining is more stable, but may fail to adapt to drift.

\subsection{Reproducibility versus Performance}

Deterministic execution improves reproducibility, but can reduce throughput or hardware efficiency.

\medskip

\noindent
Thus engineering teams often choose a balance:
\begin{itemize}
    \item maximize reproducibility in development and evaluation,
    \item allow some nondeterminism in large-scale production training when necessary.
\end{itemize}

\subsection{Automation versus Human Oversight}

Automation speeds up iteration, but some high-stakes tasks require human review, fallback rules, or manual escalation paths.

\medskip

\noindent
A strong ML system often combines:
\begin{itemize}
    \item statistical predictions,
    \item business logic,
    \item rule-based safeguards,
    \item human intervention when uncertainty is high.
\end{itemize}

\medskip

\noindent
\textbf{Debugging as a systems activity.}
Debugging in ML is broader than debugging code. One must diagnose:
\begin{itemize}
    \item data bugs,
    \item label bugs,
    \item optimization failures,
    \item evaluation leakage,
    \item serving mismatches,
    \item monitoring blind spots.
\end{itemize}

\medskip

\noindent
Useful debugging tools include:
\begin{itemize}
    \item training and validation curves,
    \item confusion matrices or error breakdowns,
    \item gradient and activation statistics,
    \item feature-distribution checks,
    \item ablation studies,
    \item inspection of individual failures.
\end{itemize}

\medskip

\noindent
\textbf{Practical debugging principle.}
When a system fails, simplify the pipeline, isolate components, verify assumptions, and test each stage independently before adding complexity back.

\medskip

\noindent
\textbf{Insight.}
Real-world ML design is governed by engineering constraints as much as by learning theory.

\section{A Unifying Perspective}

Practical machine learning systems integrate many layers of infrastructure:
\begin{itemize}
    \item \textbf{data pipelines} for collection, cleaning, labeling, and versioning,
    \item \textbf{evaluation infrastructure} for train/validation/test separation,
    \item \textbf{training pipelines} for optimization, checkpointing, and scaling,
    \item \textbf{experiment tracking} for reproducibility and comparison,
    \item \textbf{serving systems} for inference under latency and cost constraints,
    \item \textbf{monitoring systems} for drift, failures, and feedback loops.
\end{itemize}

\medskip

\noindent
\textbf{Core idea.}
A successful ML application is an end-to-end system whose components must work together reliably over time.

\medskip

\noindent
\textbf{Key insight.}
The practical effectiveness of machine learning depends not only on the model, but on how data, infrastructure, deployment, and monitoring are designed as a coherent system.

\section{Outlook}

This chapter presented machine learning as an engineering and systems discipline:
\begin{itemize}
    \item framing ML as a systems problem,
    \item building data pipelines and dataset splits,
    \item managing training workflows and experiments,
    \item designing serving systems under inference constraints,
    \item monitoring drift, failures, and feedback loops,
    \item navigating real-world engineering tradeoffs.
\end{itemize}

\medskip

This concludes the book.

\medskip

\noindent
\textbf{Final guiding principle.}
Machine learning succeeds in practice when mathematical models, data pipelines, and engineering systems are designed together rather than in isolation.

\section*{Exercises}

\begin{enumerate}
    \item \textbf{End-to-end pipeline design.}
    Choose one application:
    \begin{itemize}
        \item spam detection,
        \item medical image classification,
        \item product recommendation,
        \item fraud detection.
    \end{itemize}
    Design a complete ML pipeline for the task, including:
    \begin{itemize}
        \item data collection,
        \item cleaning and preprocessing,
        \item labeling,
        \item train/validation/test splitting,
        \item training and experiment tracking,
        \item deployment,
        \item monitoring after launch.
    \end{itemize}

    \item \textbf{Split design.}
    Explain why random splitting may be inappropriate in a time-dependent prediction task. Propose a better train/validation/test strategy.

    \item \textbf{Label quality.}
    Give three examples of labeling noise or ambiguity. For each, explain how it could affect model behavior.

    \item \textbf{Training--serving skew.}
    Describe a scenario in which the preprocessing used during serving differs from the preprocessing used in training. What failure might this create?

    \item \textbf{Inference latency.}
    A model has excellent validation accuracy but is too slow for real-time use. List several ways to reduce latency and explain the tradeoffs.

    \item \textbf{Experiment management.}
    Why is it not enough to record only final validation accuracy? What additional information should be tracked for each run?

    \item \textbf{Monitoring drift.}
    Explain the difference between covariate drift and concept drift. Give an example of each.

    \item \textbf{Feedback loops.}
    Describe how a recommendation system can create a feedback loop that changes future training data. Why can this make evaluation more difficult?

    \item \textbf{Deployment failure analysis.}
    Consider a deployed classifier whose measured accuracy suddenly drops. List at least five possible failure points in the end-to-end system, not limited to the model itself.

    \item \textbf{Debugging workflow.}
    Suppose training loss decreases, but validation performance is poor and unstable across runs. Give a systematic debugging plan that checks data, optimization, evaluation, and infrastructure.

    \item \textbf{Engineering tradeoff essay.}
    Write a short essay on one of the following tradeoffs:
    \begin{itemize}
        \item accuracy versus latency,
        \item reproducibility versus training speed,
        \item model complexity versus ease of deployment.
    \end{itemize}
    Explain why there is no universally optimal choice.

    \item \textbf{Real-world systems reflection.}
    Explain why practical ML should be viewed as a systems discipline rather than only a modeling discipline.
\end{enumerate}